
\section{Appendix D. Negative Results}

\subsection{D.1 Attempts That Failed}

\subsubsection{D.1.1 Trainable Base Weights}

\textbf{Goal.} Fine-tune base weights alongside LoRA.

\textbf{Result.} Accuracy improved initially but model diverged from pretrained 
representations; compression-like benefits lost.

\textbf{Conclusion.} Keep base frozen.

\subsubsection{D.1.2 α = 2.0 (Sparsemax)}

\textbf{Goal.} Maximize sparsity.

\textbf{Result.} Accuracy dropped to 88.30\%; training unstable.

\textbf{Conclusion.} Sparsemax is too sparse and non-smooth for stable training.

\subsubsection{D.1.3 Very Small Splash Window (W=4)}

\textbf{Goal.} Extreme local attention.

\textbf{Result.} Accuracy dropped to 75.5\%.

\textbf{Conclusion.} W=4 loses too much context.

\subsubsection{D.1.4 No Attention Sinks}

\textbf{Goal.} Simplify masking.

\textbf{Result.} Accuracy dropped to 85.5\%.

\textbf{Conclusion.} Sinks are essential for stability.

\subsubsection{D.1.5 No LoRA Adapters (r\_L=0)}

\textbf{Goal.} Test pure sparse attention.

\textbf{Result.} Accuracy dropped to 80.0\%.

\textbf{Conclusion.} LoRA recovery is essential for matching pretrained accuracy.

\subsubsection{D.1.6 Fused Bisection with Different α per Head}

\textbf{Goal.} Adaptive sparsity per head.

\textbf{Result.} No improvement over uniform α = 1.5; 2× slower.

\textbf{Conclusion.} Uniform α is sufficient.

\subsection{D.2 Abandoned Variants}

\subsubsection{D.2.1 Block-Sparse Entmax}

Tested block-structured entmax with fixed block size. Result: no improvement, 
more complex.

\subsubsection{D.2.2 Cascading Entmax}

Applied entmax multiple times in sequence. Result: accuracy collapsed; 
distribution became too sparse.

\subsubsection{D.2.3 Entmax + Softmax Mix}

Applied entmax only on some layers. Result: no clear pattern; uniform 
entmax is simpler.

\subsubsection{D.2.4 Learned α via Straight-Through}

Attempted to learn α via straight-through estimator. Result: gradient 
signal weak; no improvement.

\subsection{D.3 Failed Hyperparameters}

\subsubsection{Learning Rate}

\begin{table}[H]
\centering
\begin{tabular}{lll}
\toprule
\textbf{LR} & \textbf{DistilBERT} & \textbf{Loss stability} \\
\midrule
5e-5 & 89.5 & Stable \\
1e-4 & 90.0 & Stable \\
2e-4 & 90.25 & Stable \\
5e-4 & 89.0 & Unstable \\
1e-3 & NaN & NaN \\
\bottomrule
\end{tabular}
\end{table}

\subsubsection{Bisection Iterations}

\begin{table}[H]
\centering
\begin{tabular}{lll}
\toprule
\textbf{$n_\text{iter}$} & \textbf{Accuracy} & \textbf{Time} \\
\midrule
10 & 89.5 & Fast \\
15 & 90.10 & Medium \\
25 & 90.25 & Slow \\
50 & 90.28 & Very slow \\
\bottomrule
\end{tabular}
\end{table}

\textbf{Observation.} 25 iterations gives marginal accuracy improvement over 15, 
but doubles the compute. Consider 15 for production.

\subsubsection{Splash Radius}

\begin{table}[H]
\centering
\begin{tabular}{lll}
\toprule
\textbf{W} & \textbf{GPT-2} & \textbf{Notes} \\
\midrule
4 & 75.5 & Too local \\
8 & 77.2 & Underfit \\
16 & 78.10 & Reasonable \\
32 & 78.78 & Best \\
64 & 78.75 & Diminishing returns \\
128 & 78.80 & Full attention \\
\bottomrule
\end{tabular}
\end{table}

\subsection{D.4 Fixes That Didn't Help}

\begin{table}[H]
\centering
\begin{tabular}{lll}
\toprule
\textbf{Fix} & \textbf{Problem} & \textbf{Result} \\
\midrule
Gradient clipping & Stability & No change \\
Different LoRA init & Convergence & No change \\
Layer-wise α & Rank & No improvement \\
Position-relative splash & Flexibility & No improvement \\
Hybrid softmax+entmax & Stability & No improvement \\
\bottomrule
\end{tabular}
\end{table}

\subsection{D.5 Multi-Seed}

\begin{table}[H]
\centering
\begin{tabular}{llll}
\toprule
\textbf{Seed} & \textbf{DistilBERT} & \textbf{RoBERTa} & \textbf{GPT-2} \\
\midrule
42 & 90.25 & 77.50 & 78.78 \\
123 & 90.10 & 77.35 & 78.60 \\
456 & 90.30 & 77.60 & 78.85 \\
**Mean** & **90.22** & **77.48** & **78.74** \\
**Std** & **0.10** & **0.12** & **0.13** \\
\bottomrule
\end{tabular}
\end{table}

\textbf{Observation.} Standard deviation ~0.15\%. Results are reproducible.

\subsection{D.6 Lessons Learned}

1. \textbf{Freeze base weights} — always.
2. \textbf{α = 1.5 is optimal} — α = 2.0 is too sparse.
3. \textbf{Sinks are essential} — removing them causes 3-5\% drop.
4. \textbf{LoRA is essential} — without it, accuracy drops 5-10\%.
5. \textbf{25 bisection iterations} is a reasonable default.
6. \textbf{Not a compression method} — parameter count increases.

\subsection{D.7 Unresolved Questions}

1. \textbf{Fused CUDA kernel:} Would speedup by 10-20×?
2. \textbf{Learned α per head:} Would it help?
3. \textbf{Sparse hardware:} Real 2:4 speedup?
4. \textbf{Long context:} Does AdaSplash scale to L = 4096?
5. \textbf{Other tasks:} Generation, QA, NLI?